\section{Model architecture}\label{sec:arch}

\begin{plainwords}
All the learned drivers are built from the same small neural network, so that differences between drivers come from \emph{what they are shown} and \emph{how they were taught}, not from how big they are. This section draws the network, defines each of its inputs, and describes the two rule-based teachers.
\end{plainwords}
\begin{cautionbox}[Compared with GitHub]
The network is the teammates' \code{BEVScalarPolicy} from \pth{bev_policy/model.py} on GitHub (convolution layers of 32, 64 and 96 maps, then a layer of 128 numbers; a small network for the six vehicle-state numbers; a head of 128 units). The only difference is an optional extra-input branch added in a local, uncommitted edit; without it the network is identical to GitHub's and the two committed baseline checkpoints load unchanged. The map overlay, the collision-aware teacher and the image-only network do not exist on GitHub.
\end{cautionbox}


All learned driving policies in this project share \emph{one} network family, the teammates' \code{BEVScalarPolicy} (\pth{bev_policy/model.py}). The models differ in \emph{what they are given as input}, \emph{who labels their training data} and \emph{how they are trained},
not in the network. This is deliberate: it makes the comparison between models a comparison between input designs and training strategies. The one different network is the Image-Only Driving Model (Section~\ref{sec:catalogue}).

\subsection{The shared network}
\begin{figure}[H]\centering
\begin{tikzpicture}[font=\footnotesize,>=Stealth,
  box/.style={draw,rounded corners=2pt,minimum height=0.9cm,align=center,fill=white},
  inbox/.style={box,fill=blue!8},
  optbox/.style={box,dashed,fill=orange!8},
  outbox/.style={box,fill=green!10},
  node distance=0.45cm and 0.55cm]
\node[inbox] (bev) {Bird's-eye view\\ $1\times64\times64$};
\node[box,right=of bev] (c1) {Conv $5\times5$\\ stride 2\\ 32 maps};
\node[box,right=of c1] (c2) {Conv $3\times3$\\ stride 2\\ 64 maps};
\node[box,right=of c2] (c3) {Conv $3\times3$\\ stride 2\\ 96 maps};
\node[box,right=of c3] (gp) {Global\\ average\\ pool};
\node[box,right=of gp] (l1) {Linear\\ $96\to128$};
\node[inbox,below=1.55cm of bev] (sc) {Vehicle-state inputs\\ (3 or 6 values,\\ standardised)};
\node[box,right=of sc] (s1) {Linear\\ $n\to64$};
\node[box,right=of s1] (s2) {Linear\\ $64\to32$};
\node[optbox,below=1.45cm of s1] (ex) {Optional extra input\\ (caution, pose-error level,\\ or weather vector)};
\node[optbox,right=of ex] (e1) {Linear\\ $k\to16$};
\node[box,right=2.3cm of s2] (cat) {Concatenate\\ $128+32\,(+16)$};
\node[box,right=of cat] (h1) {Linear\\ $160\,(176)\to128$};
\node[outbox,right=of h1] (o) {Linear $128\to2$\\ tanh\\ steering,\\ throttle/brake};
\foreach \a/\b in {bev/c1,c1/c2,c2/c3,c3/gp,gp/l1,sc/s1,s1/s2,ex/e1,cat/h1,h1/o} \draw[->] (\a) -- (\b);
\draw[->] (l1.east) -- ++(0.35,0) |- (cat.north west);
\draw[->] (s2) -- (cat);
\draw[->,dashed] (e1.east) -| (cat.south);
\end{tikzpicture}
\caption{The shared policy network. Each box is one layer; every layer except the last is followed by a rectifier. The dashed boxes exist only in the models with an extra input.
Parameter count: 110\,530 without the extra input and 112\,610 with a one-dimensional extra input.}
\label{fig:network}
\end{figure}

\paragraph{Inputs.}
\begin{description}[leftmargin=3.4cm,style=nextline]
\item[Bird's-eye view] A $64\times64$ grid covering 10\,m to each side, 17.5\,m ahead and 2.5\,m behind the car (cell size 0.3125\,m). Two encodings occur (Figure~\ref{fig:inputs}):
 \emph{unmodified} (teammates' encoding: $-1$ unknown, 0 free, 0.9 ego, 1.0 detected vehicle) and \emph{map overlay} (the reserved values 0.5 for the route centre line and 0.8 for the road boundaries are filled in from the simulator map, in the ego frame; ego and detected vehicles are never overwritten).
 Vehicles are drawn as axis-aligned blocks without heading (a limitation analysed in Section~\ref{sec:failure}).
\item[Six vehicle-state inputs] speed (m/s), route progress (fraction), lateral offset to the lane centre (m), heading error (rad), curvature ahead (rad/m), distance to the goal (m); standardised with the mean and standard deviation stored in the checkpoint.
\item[Three vehicle-state inputs] The proposed design keeps only speed, route progress and distance to goal; lateral offset, heading error and curvature are \emph{not} given as numbers because the overlay already shows the road geometry relative to the car (idea borrowed from the reinforcement-learning coach of Zhang et al., 2021).
\item[Extra input] One to four additional values through the optional 16-unit branch: the caution value $c\in[0,1]$, the pose-error level(s), or a weather vector (one-hot clear/fog/rain plus intensity).
\end{description}

\explainedfigure{fig_birds_eye_view_input_examples}{0.9}{The two input encodings of the bird's-eye-view grid, drawn from one real scenario}{The model's input is a 64 by 64 grid covering 10 metres to each side, 17.5 metres ahead and 2.5 metres behind the car. Left: the teammates' encoding (unknown, free, the ego car and detected vehicles). Right: the same grid with the route centre line and road boundaries from the map added. Both are drawn from scenario seed 5169 at step 238, with the same vehicle in view.}{In the left grid the car sees free space and a vehicle ahead but nothing about the road. In the right grid it also sees where the lane edges and the planned route are.}{The map overlay gives the model the road geometry directly, which is why three vehicle-state numbers (lateral offset, heading error, curvature) could be removed from its numeric inputs. It also creates the dependence on an accurate position estimate studied above.}{fig:inputs}


\paragraph{Output.} Two numbers in $[-1,1]$ (tanh): steering, and throttle (positive) or brake (negative), the MetaDrive action contract.

\subsection{The actor--critic variant for reinforcement learning}
For reinforcement learning (\pth{bev_rl/policy.py}, class \code{BEVActorCritic}) the network above becomes the mean of a Gaussian actor: the actor's feature extractor and head are \emph{exactly} the behaviour-cloned modules, so that at initialisation the action mean equals the imitation policy's action for every input (verified in \pth{tests/test_bev_rl_policy.py}).
A state-independent learnable log standard deviation is added (initial value $-2$, i.e.\ standard deviation 0.135), and a separate critic with its own encoder (also initialised from the cloned encoder weights, so that value-loss gradients cannot disturb the actor's features) and a randomly initialised value head.

\subsection{The Image-Only Driving Model}
For the control experiment (image-only model, \pth{bev_rl/adaptive/vanilla_camera_cnn.py}) the input is only the front-left camera image (rendered at 320$\times$240, resized to 128$\times$96, scaled to $[0,1]$), processed by five convolution layers with batch normalisation (24, 36, 48, 64, 64 maps; kernels 5, 5, 3, 3, 3, all stride 2; output $64\times3\times4$) and a multilayer perceptron head with tanh output.
It receives \emph{no} bird's-eye view, speed or weather information.

\subsection{Rule-based teachers (not learned)}\label{sec:teachers}
\begin{description}[leftmargin=3.4cm,style=nextline]
\item[Stock intelligent driver model] MetaDrive's \code{IDMPolicy}: lane keeping with same-lane car following, lane changes enabled, target speed about 30\,km/h. It sees exact simulator state. It reacts only to the vehicle ahead in its own lane.
\item[Collision-aware teacher] \codename{safe_teacher} (\pth{bev_rl/adaptive/safe_teacher.py}). The stock model with lane changes \emph{disabled} plus a privileged predictive braking layer: every other vehicle within 45\,m is predicted with a constant-velocity model over a 3\,s horizon;
 if the ego rectangle (moving with its current velocity) would overlap the other vehicle's rectangle (margins 0.35\,m laterally, 1.0\,m longitudinally) at time-to-collision $t$, the throttle command is set to full brake when $t<1.3$\,s, and to a graded brake when $1.3\le t<2.4$\,s.
 It was introduced because the stock model crashes in 45\% of dense-traffic episodes, almost always into vehicles that merge, cut in or cross from the side, which a same-lane controller never reacts to (Section~\ref{sec:failure}).
\end{description}
Both teachers use information the students do not have (exact positions and velocities of all vehicles up to 45\,m, while the student's grid ends 17.5\,m ahead); this \emph{information gap} is one of the findings of the failure analysis.
